\section{Билет 27. Стоячие волны. Колебания струны с двумя закрепленными концами. Собственные частоты колебаний струны. Обертон.}

\begin{definition}
\textbf{Стоячая волна} возникает при наложении двух бегущих волн одинаковой амплитуды и частоты, распространяющихся навстречу друг другу. Уравнение стоячей волны имеет вид:
\begin{equation}
    \xi(x,t) = 2A \cos(kx) \cos(\omega t),
\end{equation}
где $A_{\text{ст}}(x) = 2A|\cos(kx)|$ --- амплитуда стоячей волны, зависящая только от координаты. Точки с нулевой амплитудой называются \textbf{узлами}, с максимальной --- \textbf{пучностями}.
\end{definition}

\begin{theorem}[Колебания струны с закрепленными концами]
Рассмотрим однородную струну длины $L$, жестко закрепленную на обоих концах ($x=0$ и $x=L$). Граничные условия требуют, чтобы смещение на концах было равно нулю для любого момента времени:
\begin{equation}
    \xi(0,t) = 0, \quad \xi(L,t) = 0.
\end{equation}
Это означает, что на концах струны всегда образуются узлы стоячей волны. Расстояние между соседними узлами равно $\lambda/2$, следовательно, на длине струны должно укладываться целое число полуволн:
\begin{equation}
    L = n \frac{\lambda_n}{2}, \quad n = 1, 2, 3, \dots \label{eq:string_condition}
\end{equation}
Отсюда получаем спектр возможных длин волн, которые могут существовать в струне:
\begin{equation}
    \lambda_n = \frac{2L}{n}. \label{eq:wavelength_string}
\end{equation}
\end{theorem}

\begin{property}[Собственные частоты и обертоны]
Частота колебаний связана с длиной волны и скоростью распространения $v$ соотношением $\nu = v/\lambda$. Подставляя \eqref{eq:wavelength_string}, находим \textbf{спектр собственных частот} струны:
\begin{equation}
    \nu_n = \frac{v}{\lambda_n} = n \frac{v}{2L}, \quad n = 1, 2, 3, \dots \label{eq:frequencies_string}
\end{equation}
\begin{itemize}
    \item При $n=1$ частота $\nu_1 = \frac{v}{2L}$ называется \textbf{основной частотой} (первой гармоникой). Это самая низкая возможная частота колебаний струны.
    \item Частоты $\nu_2, \nu_3, \dots$ называются \textbf{обертонами} (второй, третьей гармониками и т.д.). Они кратны основной частоте: $\nu_n = n \nu_1$.
\end{itemize}
Скорость волны в струне определяется натяжением $F$ и линейной плотностью $\rho$: $v = \sqrt{F/\rho}$. Следовательно, собственные частоты зависят от физических параметров струны и условий закрепления.
\end{property}

\begin{remark}
\textbf{Практическое значение:} В музыкальных инструментах звучание струны представляет собой суперпозицию основной частоты и ряда обертонов. Тембр звука определяется относительными амплитудами этих гармоник. Изменяя длину струны $L$ (зажимая пальцем), натяжение $F$ или линейную плотность $\rho$ (выбирая струны разной толщины), музыканты регулируют высоту тона. Явление дискретного спектра частот характерно не только для струн, но и для любых ограниченных волноводов (трубы, мембраны, резонаторы).
\end{remark}
